%==============================================================================
\section{The \Onehop Operator}
\label{sec:onehop}
%==============================================================================


This section presents the \onehop operator that computes the output for a 
single-hop pattern $n_i \rightarrow e \rightarrow n_j$ within an acyclic 
graph query. Specifically, given a source node $n_i(id, [prop])$ and destination 
node $n_j(id, [prop])$ connected by edge label $e$, each with its properties, 
the \onehop operator returns all tuples $(n_i, e, n_j)$, whose output size is 
bounded by $|e|$. The two nodes may be of the same or different types.
Additionally, \onehop marks output tuples that satisfy the query predicates 
$p_s$, $p_e$, and/or $p_d$ with a flag, leaving non-matching tuples unmarked. 
This flagging mechanism is central to hiding intermediate result sizes, as we 
describe in []. The \onehop operator reveals no information 
beyond the sizes of the source node table, destination node table, and the 
edge list --- all of which are public in our threat model.


Recall that similar to plaintext \gdbs~\cite{janusgraph, feng2023kuzu, 
deutsch2019tigergraph, sun2015sqlgraph}, \system also stores two versions of each 
edge list but without revealing the graph structure: a \textit{forward list} 
$E_F$ sorted by source node ids and a \textit{reverse list} $E_R$ sorted by 
destination node ids, both of size $|e|$.
The \onehop operator exploits both lists to parallelize one-hop query 
processing. We first describe a plaintext algorithm and then identify its 
leakages, which motivate our oblivious design.


\noindent\textbf{Plaintext algorithm.}
A natural plaintext approach proceeds as follows. Inspired by hash joins, 
store source node entries in a hash table and probe it once per entry in 
$E_F$, attaching matching source properties to each edge; the destination side
is handled symmetrically against $E_R$ in parallel. Then, sort one of the lists,
say $E_R$ on source ids, and merge the two lists row-wise, combining source node,
edge, and destination node properties into the output tuple. Finally, apply 
predicate filters on the source, destination, and edge properties, retaining 
only tuples that pass. This enables bidirectional parallel processing of the 
traversal.

\noindent\textbf{Leakages.}
The plaintext algorithm above has three sources of leakage that make it 
non-oblivious:
\begin{enumerate}[leftmargin=*]
    \item \textit{Hash table access pattern.} Hash table construction and probing
    reveals for each probe whether a match occurred and, if so, at which location,
    exposing which node ids appear in the edge list.

    \item \textit{Node degree.} A node id appears in the edge list once per 
    incident edge, so high-degree nodes are probed repeatedly. The number 
    of repeated probes to the same id reveals the degree distribution and 
    thus the graph structure.

    \item \textit{Predicate selectivity.} Filtering the merged output on 
    node and edge predicates reveals how many tuples satisfy predicates, 
    leaking the selectivity of query predicate values.
\end{enumerate}

We design the \onehop operator to retain the bidirectional parallelism of 
the plaintext algorithm while provably eliminating all three leakages.


% This
% creates two problems. First, the oblivious hash table we build [refer this] on permits each
% key to be queried \emph{at most once}, so repeated probes for the same source
% would violate its security precondition. Second, even setting that aside, the
% number of probes for a given id would equal that id's degree, directly exposing
% the graph's degree distribution. The operator resolves both by probing each
% distinct key \emph{exactly once} and then obliviously propagating the retrieved
% attributes to the remaining edges that share that key, so the access pattern
% depends only on $|E|$ --- never on any node's degree.


% An oblivious \onehop must hide three sources of leakage. First,
% \emph{hash-bucket access}: the memory pattern of retrieving a node's attributes
% by id would expose which key was probed. Second, \emph{node degree}: as pointed above, the number
% of source-side (or destination-side) lookups for a given id would equal its
% degree, leaking the degree distribution of the graph. Third, \emph{predicate
% selectivity}: the number of edges that satisfy the query predicates would reveal
% how many tuples match. The algorithm below neutralizes the first two by
% construction; the third is handled by always emitting $|E|$ rows and deferring
% any selectivity-dependent sizing, as we discuss at the end of the section.

%------------------------------------------------------------------------------
\subsection{\Onehop Algorithm}
\label{sec:onehop:algorithm}


We address the three leakages in order and present the operator in Algorithm~\ref{alg:onehop}.

\noindent\textbf{Addressing leakage 1: Hash table access pattern.}
Our algorithm stores source and destination node information in oblivious 
hash tables (\ohts). Our implementation utilizes H2O2RAM~\cite{zheng2025h2o2ram}
but the operator can use any \oht.
\ohts ensure that both table construction and probing are oblivious: the 
access pattern depends only on the input size, not on the values stored, and 
probing reveals neither whether a match was found nor the location of any 
matched entry. Note that \oht construction requires unique input keys, which is satisfied 
here since node ids used as keys for both source and destination tables
are unique by definition. However, \ohts do not protect against \emph{repeated probes 
of the same key} --- if the same node id is probed multiple times due to its 
degree in the edge list, that repetition remains visible. Node degree leakage 
therefore persists and is addressed next. 

\noindent\textit{Offline/online separation.}
\oht construction is generally expensive. For example, H2O2RAM's build complexity 
ranges from $O(n\log^2 n)$ for small tables to 
$O(n\log^2\log n)$ for large ones~\cite{zheng2025h2o2ram}. We observe, 
however, that in our operator the hash table acts as a \emph{index}
built on node entries and then probed by edge entries at query time. 
This means index construction can occur \emph{offline}, before any query 
arrives. Furthermore, once an index is built for nodes $n_i$ and $n_j$ 
paired with edge label $e$, it can be reused across all queries over the 
same node-edge combination, amortizing the construction cost. Note that 
the same \oht cannot be reused across different edge labels without 
reconstruction, so one index is required per $(n, e)$ pair. Changes to
a node or an edge also requires a reconstruction.




\begin{algorithm*}[t]
\caption{Oblivious \onehop Operator}
\label{alg:onehop}
\begin{algorithmic}[1]
\Require Forward list $E_{\!f}$ (sorted by $src\_id$), reverse list $E_{\!r}$
         (sorted by $dst\_id$); precomputed indices $H_S$, $H_D$; optional
         predicates $P$
\Ensure  Result $R$ with $|R| = |E|$
\Statex
\State \textbf{Source branch} \Comment{concurrent with the destination branch}
\State $\textsc{Dedup}(E_{\!f})$: for each row $i$,\algoIndent{1}
       $E_{\!f}[i].k \gets \textsc{OCmpSet}\big(E_{\!f}[i].src\_id = E_{\!f}[i\!-\!1].src\_id,\ \bot,\ E_{\!f}[i].src\_id\big)$
\State for each row $i$: probe $H_S[E_{\!f}[i].k]$ \Comment{all $|E|$ rows; dummy keys are no-ops}
\State $\textsc{Redup}(E_{\!f})$: for each row $i$,\algoIndent{1}
       $E_{\!f}[i].sattr \gets \textsc{OCmpSet}\big(E_{\!f}[i].k = \bot,\ E_{\!f}[i\!-\!1].sattr,\ E_{\!f}[i].sattr\big)$
\State $\textsc{Merge}(E_{\!f}, \text{source attributes})$ \Comment{attach source columns}
\Statex
\State \textbf{Destination branch} \Comment{symmetric over $dst\_id$, $H_D$}
\State $\textsc{Dedup}$ / probe $H_D$ / $\textsc{Redup}$ / $\textsc{Merge}$ on $E_{\!r}$
\State $E_{\!r} \gets \textsc{OSort}(E_{\!r})$ into edge order \Comment{re-align to source branch}
\Statex
\State \textbf{Combine}
\State $R \gets \textsc{Merge}(E_{\!f}, E_{\!r})$ \Comment{column-wise, aligned per edge}
\If{$P \neq \emptyset$} \Comment{standalone query only; skipped under decomposition}
    \State $R \gets \textsc{OCompact}\big(\text{rows of } R \text{ satisfying } P\big)$
\EndIf
\State \Return $R$
\end{algorithmic}
\end{algorithm*}

\noindent\textbf{Addressing leakage 2: Node degree.}
Since \ohts cannot protect against repeated probes to the same node id, the 
\onehop operator addresses node degree leakage as the first step of online 
processing. The key idea is to \emph{deduplicate} $E_F$ (and $E_R$ in 
parallel), replacing all but the first occurrence of each node id with a 
dummy, so that each distinct node id is probed exactly once.

Deduplication exploits the fact that $E_F$ is already sorted by source id 
and $E_R$ by destination id. A linear scan over each list compares each 
entry with its predecessor and replaces repeated ids with dummies using 
\cmpset, leaving the first occurrence of each distinct id as the only real
entry per run.

We then probe the hash table index \emph{once per row}, including dummies.
A dummy key 
resolves to a no-op that the table absorbs, while a real key retrieves its 
node attributes and appends them to the edge attributes. Probing every row, 
rather than only real ones, fixes the 
total probe count at $|E|$ regardless of the degree distribution, necessary
for maintaining obliviousness. 
\emph{Reduplication} then restores suppressed rows. Specifically,
a second linear scan 
propagates, via \cmpset, the retrieved attributes from each real row forward 
into all subsequent dummy rows in the same run, so every edge row ends up 
annotated with its source (or destination) node attributes. 

The source-side branch produces an annotated edge list ordered by $src\_id$, 
while the destination-side branch produces one ordered by $dst\_id$. To 
align them for merging, we obliviously sort the destination list by $src\_id$. 
A final \textsc{Merge} then combines the two lists row-wise, attaching both 
source and destination node attributes to each edge entry. This single 
re-alignment sort is the only oblivious sort in the entire \onehop operator.

\noindent\textbf{Addressing leakage 3: Predicate selectivity.}
To hide predicate selectivity, the operator linearly scans the merged list 
and annotates each row with a boolean flag indicating whether it satisfies 
the query predicates. What happens next depends on the operator's role in 
the query:

\begin{itemize}[leftmargin=*]
    \item {Standalone query.} The algorithm obliviously compacts the 
    flagged list and returns the result to the user. The output size equals 
    the final query result size, which is public under our threat model.

    \item {Component of a larger query.} Compaction is \emph{not} 
    performed. The output remains exactly $|E|$ rows, keeping intermediate sizes 
    independent of predicate 
    selectivity. The boolean flags are preserved for the downstream 
    multi-way join, which uses them to hide the final output size, as 
    discussed in [].
\end{itemize}

In both cases, the adversary learns nothing about how many rows satisfy the 
predicates until the final result size is deliberately released. \hfill$\Box$



% \paragraph{Two branches over the pre-sorted lists.}
% Because the forward and reverse edge lists are already maintained in sorted order
% (\Cref{sec:background:graph}), the operator needs no online sort to group edges by
% key --- the grouping comes for free from the storage layout. The source branch
% consumes the forward list and the destination branch the reverse list, and the
% two run concurrently on independent working copies.

% \paragraph{Deduplicate, probe, reduplicate.}
% Within a branch, \emph{deduplication} enforces the single-query-per-key
% constraint: in a linear scan over the sorted list, we compare each row's key with
% its predecessor and, using \textsc{OCmpSet}, obliviously replace it with a dummy
% ($\bot$) whenever the two match. Exactly one row per distinct key, the first in
% its run, keeps a live key. We then probe the index \emph{once for every row},
% dummies included; a dummy key resolves to a no-op the table absorbs, while live
% keys retrieve their node attributes. Probing every row, rather than only the live
% ones, keeps the query count equal to $|E|$ regardless of the graph's degree
% distribution; probing only live keys would let the query count leak the number of
% distinct sources. \emph{Reduplication} then refills the suppressed rows: a second
% scan copies, via \textsc{OCmpSet}, the probed attributes from the preceding row
% whenever the current key is a dummy, propagating each run's attributes forward. A
% \textsc{Merge} attaches the retrieved node attributes to the edge rows as new
% columns.

% \paragraph{Align and combine.}
% The source branch finishes in $src\_id$ order, but the destination branch finishes
% in $dst\_id$ order; we therefore oblivious-sort the destination branch back into
% the same per-edge order so the two $|E|$-row tables align row-for-row. A final
% \textsc{Merge} then combines them column-wise, giving every edge both its source
% and destination attributes. This single re-alignment is the operator's only online
% sort.

% \paragraph{Filtering.}
% When the operator answers a standalone one-hop query, we evaluate the predicates
% $P$ over every row after the combine and obliviously compact the matches; the
% result size is the query's (public) output size. When the operator instead serves
% as a building block under query decomposition (\Cref{sec:decomposition}), no
% predicates are applied here, and hence, the output remains exactly $|E|$ rows, keeping
% intermediate sizes independent of selectivity.

%------------------------------------------------------------------------------
% Alternatives considered (closes the section).
%------------------------------------------------------------------------------
The one-hop pattern could in principle be evaluated as two sequential 
oblivious binary joins using an operator such as 
Obliviator~\cite{mavrogiannakis2025obliviator}, at equal security but with 
two full join executions and an intermediate materialization. Our operator 
replaces this with one pass of probes per side and a single re-alignment 
sort, yielding a $1.1$--$1.7\times$ latency improvement over the two-join 
baseline (\Cref{sec:evaluation}).

